\begin{lemma}
Let $\Omega\subseteq\mathbb{C}$ be open, let $u\colon\mathbb{C}\to\mathbb{R}$
and let $F_0,F_1\colon\mathbb{C}\to\mathbb{C}$.  Assume
\begin{enumerate}
\item[(H1)] $u\in L^1_{\mathrm{loc}}(\Omega)$ in the sense of
  \noderef{MemLpLocOn};
\item[(H2)] $F_0$ and $F_1$ are continuous on $\Omega$;
\item[(H3)] for $j=0,1$, $F_j$ is the distributional $j$-th partial derivative
  of $u$ on $\Omega$, in the sense of \noderef{IsWeakPartialOn}.
\end{enumerate}
Then there is $w\colon\mathbb{C}\to\mathbb{R}$ with
\begin{enumerate}
\item[(a)] $w(z)=u(z)$ for almost every $z$ with $z\in\Omega$;
\item[(b)] $w$ of class $C^1$ on $\Omega$;
\item[(c)] $\partial_jw=F_j$ on $\Omega$ for $j=0,1$, where $\partial_0$ and
  $\partial_1$ are the partial derivatives of \noderef{PartialWord} along the
  words $[0]$ and $[1]$, and where the real number $\partial_jw(z)$ is read as
  a complex number.  Clause (c) in particular asserts that $F_0$ and $F_1$ are
  real-valued on $\Omega$; that is a consequence of (H3), proved as Step 0
  below, and not an extra assumption.
\end{enumerate}

This is the first of the three steps Section 6 of the paper states in one line.
It is what makes the conclusion of Theorem 1 a true statement about $u$: a
$p$-harmonic function is constrained only through integrals, so it may be
altered on a null set without ceasing to be $p$-harmonic, while membership in a
H\"older class is not invariant under such an alteration; the regularity
assertion must therefore be made of a representative.  For a $p$-harmonic $u$
the hypothesis (H2) holds because the complex gradient $f$ is continuous, so
that $\nabla u=(2\operatorname{Re}f,-2\operatorname{Im}f)$ is continuous on all
of $\Omega$.  The instance $F_0=F_1=0$ is the constant-$u$ branch.
\end{lemma}
\begin{proof}
Write $d\sigma$ for plane Lebesgue measure, $B(z,r)=\{\zeta:|\zeta-z|<r\}$ and
$\bar B(z,r)=\{\zeta:|\zeta-z|\le r\}$.

\emph{Step 0: $F_0$ and $F_1$ are real-valued on $\Omega$.}  Let $j\in\{0,1\}$
and let $\varphi$ be a test function for $\Omega$ in the sense of
\noderef{IsTestFn} which is real-valued.  Then $\partial_j\varphi$ is
real-valued, so $\int u\,\partial_j\varphi\,d\sigma$ is a real number; by (H3)
it equals $-\int F_j\varphi\,d\sigma$, whence
$\int(\operatorname{Im}F_j)\varphi\,d\sigma=0$.  Now
$\operatorname{Im}F_j$ is continuous on $\Omega$ by (H2); if it were nonzero at
some $z\in\Omega$, say positive, it would be positive on a ball
$B(z,r)\subseteq\Omega$, and taking for $\varphi$ a nonnegative real test
function supported in that ball and positive at $z$ --- a standard smooth bump
--- would make $\int(\operatorname{Im}F_j)\varphi\,d\sigma>0$, a
contradiction; the negative case is symmetric.  Hence
\begin{equation}\tag{0}
  F_j(z)\in\mathbb{R}\qquad(z\in\Omega,\ j=0,1).
\end{equation}

\emph{Step 1: the local set-up.}  Fix $z_*\in\Omega$.  Since $\Omega$ is open
there is $\rho>0$ with $\bar B(z_*,4\rho)\subseteq\Omega$.  Put
\[
  U=B(z_*,2\rho),\qquad K=\bar B(z_*,2\rho),
\]
so $K$ is compact and convex, $K\subseteq B(z_*,3\rho)$, and
$\bar B(z_*,3\rho)\subseteq\Omega$.  Let
\[
  \hat u=\mathbf{1}_{\bar B(z_*,4\rho)}\,u,\qquad
  \hat F_j=\mathbf{1}_{\bar B(z_*,4\rho)}\,F_j .
\]
By (H1) and compactness of $\bar B(z_*,4\rho)\subseteq\Omega$ the function
$\hat u$ is integrable on $\mathbb{C}$, hence locally integrable; and $\hat F_j$
is bounded and measurable, hence also locally integrable, $F_j$ being
continuous on the compact $\bar B(z_*,4\rho)$.

For $0<\varepsilon<\rho$ let $\rho_\varepsilon\colon\mathbb{C}\to\mathbb{R}$ be
a smooth mollifier: $\rho_\varepsilon\ge0$, smooth, supported in
$\bar B(0,\varepsilon)$, with $\int\rho_\varepsilon\,d\sigma=1$ (the normalized
smooth bump functions).  Put
\[
  u_\varepsilon=\rho_\varepsilon*\hat u,\qquad
  G_{j,\varepsilon}=\rho_\varepsilon*\hat F_j .
\]
Both are $C^\infty$ on $\mathbb{C}$, being convolutions of a smooth compactly
supported function with a locally integrable one, and $u_\varepsilon$ is
real-valued because $\hat u$ and $\rho_\varepsilon$ are.  Note that for
$x\in B(z_*,3\rho)$ and $\varepsilon<\rho$ the convolution integrals see only
points of $\bar B(x,\varepsilon)\subseteq\bar B(z_*,4\rho)$, where $\hat u=u$
and $\hat F_j=F_j$.

\emph{Step 2: $\partial_ju_\varepsilon=G_{j,\varepsilon}$ on $B(z_*,3\rho)$.}
Let $x\in B(z_*,3\rho)$ and put $\varphi_x(z)=\rho_\varepsilon(x-z)$.  Then
$\varphi_x$ is smooth with support contained in
$\bar B(x,\varepsilon)\subseteq\bar B(z_*,4\rho)\subseteq\Omega$, since
$|x-z_*|<3\rho$ and $\varepsilon<\rho$; so $\varphi_x$ is a test function for
$\Omega$ in the sense of \noderef{IsTestFn}.

Differentiating the convolution under the integral sign --- legitimate because
$\rho_\varepsilon$ is $C^\infty$ with compact support and $\hat u$ is
integrable --- gives
\[
  \partial_ju_\varepsilon(x)
  =\int\bigl(\partial_j\rho_\varepsilon\bigr)(x-z)\,\hat u(z)\,d\sigma(z).
\]
On the other hand $\partial_{z_j}\varphi_x(z)=-(\partial_j\rho_\varepsilon)(x-z)$
by the chain rule, so, using that $\varphi_x$ and $\partial_j\varphi_x$ are
supported where $\hat u=u$ and $\hat F_j=F_j$, and then (H3),
\[
  \partial_ju_\varepsilon(x)
  =-\int u\,\partial_j\varphi_x\,d\sigma
  =\int F_j\,\varphi_x\,d\sigma
  =\int\hat F_j(z)\,\rho_\varepsilon(x-z)\,d\sigma(z)
  =G_{j,\varepsilon}(x).
\]
Thus
\begin{equation}\tag{1}
  \partial_ju_\varepsilon=G_{j,\varepsilon}\quad\text{on }B(z_*,3\rho),
  \qquad j=0,1,\ 0<\varepsilon<\rho .
\end{equation}

\emph{Step 3: $G_{j,\varepsilon}\to F_j$ uniformly on $K$.}  The function $F_j$
is continuous on the compact set $\bar B(z_*,3\rho)$, hence uniformly
continuous there; let $\omega$ be a modulus of continuity, so
$|F_j(a)-F_j(b)|\le\omega(|a-b|)$ for $a,b\in\bar B(z_*,3\rho)$ with
$\omega(t)\to0$ as $t\to0^{+}$.  For $x\in K$ and $0<\varepsilon<\rho$, using
$\int\rho_\varepsilon=1$, $\rho_\varepsilon\ge0$ and
$\operatorname{supp}\rho_\varepsilon\subseteq\bar B(0,\varepsilon)$,
\[
  \bigl|G_{j,\varepsilon}(x)-F_j(x)\bigr|
  =\Bigl|\int\rho_\varepsilon(y)\bigl(F_j(x-y)-F_j(x)\bigr)d\sigma(y)\Bigr|
  \le\sup_{|y|\le\varepsilon}\bigl|F_j(x-y)-F_j(x)\bigr|\le\omega(\varepsilon),
\]
the points $x-y$ lying in $\bar B(z_*,3\rho)$ because $|x-z_*|\le2\rho$ and
$|y|\le\varepsilon<\rho$.  So
\begin{equation}\tag{2}
  \sup_{K}\bigl|G_{j,\varepsilon}-F_j\bigr|\le\omega(\varepsilon)
  \ \longrightarrow\ 0\qquad(\varepsilon\to0^{+}).
\end{equation}

\emph{Step 4: $u_\varepsilon\to u$ almost everywhere on $U$.}  The
mollifications of a locally integrable function converge to it almost
everywhere as the radius tends to $0$.  Applied to $\hat u$ this gives a null
set $N$ with
\begin{equation}\tag{3}
  u_{\varepsilon}(x)\to\hat u(x)=u(x)\qquad(\varepsilon\to0^{+}),
  \quad x\in U-N ,
\end{equation}
the identification $\hat u=u$ holding on $U\subseteq\bar B(z_*,4\rho)$.

\emph{Step 5: the local representative.}  Fix a sequence
$\varepsilon_n=\rho/(n+2)\downarrow0$ and choose a base point $z_1\in U-N$;
this is possible because $U$ has positive measure and $N$ is null.  Let
$h=u_{\varepsilon_n}-u_{\varepsilon_m}$.  By (1), $h$ is $C^\infty$ on
$B(z_*,3\rho)$ with $\partial_jh=G_{j,\varepsilon_n}-G_{j,\varepsilon_m}$
there, so by (2), writing
$\sigma_{n,m}=2\bigl(\omega(\varepsilon_n)+\omega(\varepsilon_m)\bigr)$,
\[
  \sup_{K}\bigl\|\mathrm{d}h\bigr\|
  \le\sup_{K}\bigl|\partial_0h\bigr|+\sup_{K}\bigl|\partial_1h\bigr|
  \le\sigma_{n,m},
\]
since an $\mathbb{R}$-linear map $\mathbb{C}\to\mathbb{C}$ has operator norm at
most the sum of the moduli of its values on $1$ and $i$.  The set $K$ is convex
and $h$ is differentiable at each of its points, so the mean value inequality
gives, for $x\in K$,
\[
  \bigl|h(x)-h(z_1)\bigr|\le\sigma_{n,m}\,|x-z_1|\le4\rho\,\sigma_{n,m}.
\]
Moreover $h(z_1)=u_{\varepsilon_n}(z_1)-u_{\varepsilon_m}(z_1)\to0$ as
$n,m\to\infty$ by (3), since $z_1\in U-N$.  As $\sigma_{n,m}\to0$ as
$n,m\to\infty$, the sequence $(u_{\varepsilon_n})$ is uniformly Cauchy on $K$,
hence converges uniformly on $K$ to a continuous real-valued function $w_*$.

Three consequences.  First, by (3), for $x\in U-N$ we have
$u_{\varepsilon_n}(x)\to u(x)$ and $u_{\varepsilon_n}(x)\to w_*(x)$, so
\begin{equation}\tag{4}
  w_*=u\qquad\text{almost everywhere on }U .
\end{equation}
Second, let $L(x)$ denote the $\mathbb{R}$-linear map $\mathbb{C}\to\mathbb{C}$
determined by $L(x)[1]=F_0(x)$ and $L(x)[i]=F_1(x)$, which is real-valued by
(0).  By (1) each $u_{\varepsilon_n}$ is differentiable at each point of the
open set $U$ with derivative the map sending $1,i$ to
$G_{0,\varepsilon_n},G_{1,\varepsilon_n}$; by (2) those derivatives converge to
$L$ uniformly on $K\supseteq U$, hence locally uniformly on $U$; and
$u_{\varepsilon_n}\to w_*$ pointwise on $U$.  By the theorem that a locally
uniform limit of derivatives is the derivative of the limit, $w_*$ is
differentiable at every $x\in U$ with $\mathrm{d}w_*(x)=L(x)$, that is
\begin{equation}\tag{5}
  \partial_jw_*=F_j\qquad\text{on }U,\ j=0,1 .
\end{equation}
Third, by (5) and (H2) the derivative $\mathrm{d}w_*$ is continuous on $U$, so
\begin{equation}\tag{6}
  w_*\ \text{is of class }C^1\text{ on }U .
\end{equation}

\emph{Step 6: a canonical global definition.}  Define
\[
  w(z)=\lim_{r\to0^{+}}\frac{1}{\pi r^{2}}\int_{B(z,r)}u\,d\sigma
  \quad(z\in\Omega),\qquad w(z)=0\quad(z\notin\Omega),
\]
and let us check that the limit exists for every $z\in\Omega$ and that $w$
agrees with every local representative.

Let $z_*\in\Omega$ and let $U$, $w_*$ be as in Steps 1--5 for this $z_*$.  Let
$z\in U$.  For $r>0$ small enough that $B(z,r)\subseteq U$, (4) gives
\[
  \frac{1}{\pi r^{2}}\int_{B(z,r)}u\,d\sigma
  =\frac{1}{\pi r^{2}}\int_{B(z,r)}w_*\,d\sigma ,
\]
the two integrands agreeing almost everywhere on $B(z,r)\subseteq U$; and since
$w_*$ is continuous at $z$ by (6), the averages of $w_*$ over $B(z,r)$ tend to
$w_*(z)$ as $r\to0^{+}$.  Hence the limit defining $w(z)$ exists and
\begin{equation}\tag{7}
  w=w_*\qquad\text{on }U .
\end{equation}
In particular $w$ is real-valued, since $w_*$ is.

\emph{Step 7: conclusion.}  Each $z_*\in\Omega$ has a neighbourhood
$U(z_*)=B(z_*,2\rho(z_*))$ as in Step 1, and by (7), (6), (5),
\[
  w=w_*^{(z_*)}\ \text{on }U(z_*),\qquad
  w\in C^1\bigl(U(z_*)\bigr),\qquad \partial_jw=F_j\ \text{on }U(z_*).
\]
Since $\Omega$ is open and is covered by these neighbourhoods, $w$ is of class
$C^1$ on a neighbourhood of each of its points in $\Omega$, hence $C^1$ on
$\Omega$, which is (b); and $\partial_jw=F_j$ at every point of $\Omega$, which
is (c).

For (a): $\Omega$ is an open subset of $\mathbb{C}$, hence second countable, so
the cover $\{U(z_*)\}_{z_*\in\Omega}$ has a countable subcover
$U_1,U_2,\dots$ of $\Omega$.  By (4) and (7), for each $i$ there is a null set
$N_i$ with $w=u$ on $U_i-N_i$.  The set $N=\bigcup_iN_i$ is null, being a
countable union of null sets, and every $z\in\Omega-N$ lies in some $U_i$ and
hence satisfies $w(z)=u(z)$.  So $w=u$ almost everywhere on $\Omega$, which is
(a).
\end{proof}
